\chapter{Orchestration et CI/CD}

\section{Airflow}
Airflow représente le pipeline sous forme de DAG : chaque tâche a un état, des logs et des dépendances \cite{airflow_tasks}. Cette vue est importante pour le rapport car elle prouve que le pipeline peut être orchestré, pas seulement exécuté à la main.

\screenshotfigure[0.62\textheight]{14_airflow_dag_success.png}{DAG Airflow exécuté avec succès.}{fig:airflow_success}
\screenshotfigure[0.62\textheight]{21_airflow_dag_graph_view.png}{Vue graphe Airflow avec dépendances.}{fig:airflow_graph}

\begin{lstlisting}
docker compose --profile airflow up -d airflow
docker exec cyberguard-airflow bash -lc "airflow dags list"
docker exec cyberguard-airflow bash -lc "airflow dags test cyberguard_mlops_pipeline 2026-05-06"
\end{lstlisting}

\section{GitHub Actions}
GitHub Actions lance un workflow à chaque push, pull request ou déclenchement manuel \cite{github_actions_syntax}. Le workflow vérifie le formatage, le linting, le typage, les tests, la génération d'artefacts et les builds Docker.

\screenshotfigure[0.58\textheight]{15_github_repo_actions.png}{GitHub Actions avec workflow réussi.}{fig:github_actions}
\smallshot{15_github_repo.png}{Dépôt GitHub CyberGuard MLOps.}{fig:github_repo}

\begin{table}[H]
\centering
\caption{Étapes GitHub Actions}
\label{tab:ci_steps}
\begin{tabularx}{\linewidth}{@{}lY@{}}
\toprule
\textbf{Étape} & \textbf{Description}\\
\midrule
Install & Installe Python 3.11 et les dépendances.\\
Format/lint & Black, isort, Ruff, MyPy.\\
Artifacts & Ingestion, validation, entraînement.\\
Tests & Pytest.\\
Build API & Construction image Docker FastAPI.\\
Build SOC UI & Construction image Docker Streamlit.\\
\bottomrule
\end{tabularx}
\end{table}

\begin{lstlisting}
gh workflow run ci.yml --repo Fadi-AICH/cyberguard-mlops
gh run list --repo Fadi-AICH/cyberguard-mlops --limit 5
\end{lstlisting}
